%%%%%%%%%%%%%%%%%%%%%%%%%%%%%%%%%%%%%%%%%%%%%%%%%
\section{Coordinate-space quasi-PDFs}
\label{sec:def}

The interest and excitement of studying quasi-PDFs is based on the potential of extracting PDFs and hadron structure from the first principle calculation of LQCD.  Consequently, the value for studying quasi-PDFs depends on the reliability and accuracy of the matching between quasi-PDFs and PDFs, as proposed in Eq.~(11) of Ref.~\cite{Ji:2013dva},
\begin{equation}
\label{eq:matching_ji}
\tilde{q}(x,\mu^2,P^z) = \int_x^1 \frac{dy}{y} Z\left(\frac{x}{y},\frac{\mu}{P^z}\right) q(y,\mu^2) \, ,
\end{equation}
with $Z(x,\mu/P^z)=\delta(x-1)+(\alpha_s/2\pi)Z^{(1)}(x,\mu/P^z) +\dots \to \delta(x-1)$ as $P^z\to \infty$ for quark distribution.  This is equivalent to require the reliability and accuracy of the PQCD factorization of the quasi-PDFs,
\begin{eqnarray}
\label{eq:factorization}
\tilde{f}_{i/p}(\tilde{x},\tilde{\mu}^2,p_z) & \approx &
\sum_{j} \int_0^1 \frac{dx}{x} {\cal C}_{ij}(\frac{\tilde{x}}{x},\tilde{\mu}^2,\mu^2,p_z)
\nonumber \\
&& \hskip 0.3in
\times f_{j/p}(x,\mu^2) + {\cal O}\left(\frac{1}{\tilde{\mu}^2}\right),
\end{eqnarray}
where $i, j =q, \bar{q}, g$, ${\cal C}$'s are IR safe and perturbatively calculated matching coefficients,  $\tilde{\mu}$ is the renormalization scale of quasi-PDFs, $\mu (\sim{\hskip -0.05in}\tilde{\mu})$ is the factorization scale of PDFs, and the power correction, ${\cal O}(1/\tilde{\mu}^2)$, is characterized by the size of high-twist quark-gluon correlation functions, like in all collinear PQCD factorization formalisms. We found that if the perturbative UV divergences of quasi-PDFs in the left-hand-side (LHS) of Eq.~(\ref{eq:factorization}) are regularized by a momentum cutoff, the factorized convolution in the right-hand-side (RHS) is well behaved. However, the momentum cutoff regulator is hard to implement consistently at higher order calculations in perturbative expansion.  On the other hand, we noticed that if we regularize the perturbative UV divergences of quasi-PDFs by dimensional regularization (DR), the integration over the momentum fraction $x$ on the RHS of the Eq.~\eqref{eq:factorization} will be divergent. This is because sea-quark and gluon PDFs, $f_{j/p}(x,\mu^2) \to x^{-\alpha}$ as $x\to0$, with $1<\alpha<2$, and the calculated matching coefficient, ${\cal C}^{(1)}_{ij}(\tilde{x}/x,\tilde{\mu}^2,\mu^2,p_z)$, evaluated in using DR, has a term proportional to $1/(\tilde{x}/x)$ as $x\to 0$  \cite{Ma:2014jla,Xiong:2013bka}.  As a result, the lower end of the $x$-integration in Eq.~(\ref{eq:factorization}) leads to the divergence, $\int_0 dx/x\, (x/\tilde{x})\, x^{-\alpha} \to \infty$.

By applying the factorization formula in Eq.~(\ref{eq:factorization}) to an asymptotic parton state, ``$j$'', and expanding the both sides of the factorized equation to the first order in power of $\alpha_s$, we have ${\cal C}^{(1)}_{ij}(y,\tilde{\mu}^2,\mu^2,p_z)=\tilde{f}^{(1)}_{i/j}(y,\tilde{\mu}^2,p_z) - f^{(1)}_{i/j}(y,\mu^2)$. As $f^{(1)}_{i/j}(y,\mu^2)$ vanishes as $y\to\infty$, the asymptotic behavior of ${\cal C}^{(1)}_{ij}(\tilde{x}/x,\tilde{\mu}^2,\mu^2,p_z)$ as $x\to 0$ is fully determined by the large $\tilde{x}$ behavior of $\tilde{f}^{(1)}_{i/j}(\tilde{x},\tilde{\mu}^2,p_z)$.

As $\tilde{x}\to\infty$, the large momentum behavior of quasi-PDFs, defined in Eqs.~\eqref{eq:ftq} and \eqref{eq:ftg}, is closely related to the asymptotic behavior of the hadronic matrix elements as the separation of the two fields, $\xi_z\to 0$.  The divergence in the $x$-convolution of the factorized formalism in Eq.~(\ref{eq:factorization}), found above, indicates that hadronic matrix elements could be ill-defined when $\xi_z\to 0$, which is indeed the case as we will demonstrate in the next section.  As pointed out in Ref.~\cite{Ma:2014jla}, the coordinate-space quasi-PDFs are better quantities for the calculation in LQCD, for the discussion of renormalization, and for the extraction of PDFs.

We define coordinate-space quasi-PDFs as following,
\begin{align}\label{eq:ftqx}
\begin{split}
&\tilde{F}_{q/p}(\xi_z,\tilde{\mu}^2,p_z)\\
=& \frac{e^{i p_z \xi_z}}{p_z} \langle h(p)| \overline{\psi}_q(\xi_z)\,\frac{\gamma_z}{2}
\Phi_{n_z}^{(f)}(\{\xi_z,0\})\, \psi_q(0) | h(p) \rangle,
\end{split}
\end{align}
and
\begin{align}\label{eq:ftgx}
\begin{split}
&\tilde{F}_{g/p}(\xi_z,\tilde{\mu}^2,p_z)\\
=&\frac{e^{i p_z \xi_z}}{p_z^2}\langle h(p)| F_{z}^{\ \nu}(\xi_z)\,
\Phi_{n_z}^{(a)}(\{\xi_z,0\})\, F_{z\nu}(0) | h(p) \rangle,
\end{split}
\end{align}
where $\tilde{\mu}$ is a renormalization scale, and $\Phi_{n_z}^{(f,a)}(\{\xi_z,0\})$ are gauge links, defined in Sec.~\ref{sec:intro}. We also define $n_z^\mu = (0,0_\perp, 1)$ with $g^{\mu\nu}=\text{diag}(1,-1,-1,-1)$, and $v\cdot n_z =- v_z$ for any vector $v^\mu$, and we have $n_z^2 = -1$.


We will demonstrate in the next section that coordinate-space quasi-PDFs are well defined for finite $\xi_z$, while they are divergent when $\xi_z\to0$.  if one wants to have a well defined momentum-space quasi-PDFs, one has to properly and consistently subtract off the divergence at $\xi_z\to0$.  From the definitions of quasi-PDFs, we have
\begin{align}
\tilde{F}_{i/p}(-\xi_z,\tilde{\mu}^2,p_z)=\tilde{F}_{i/p}^*(\xi_z,\tilde{\mu}^2,p_z).
\end{align}
To take the advantage of this relation and simplify the discussion, we assume $\xi_z>0$ in the following calculation.  For the final result, we will then express it in the form that is correct for arbitrary value of $\xi_z$.

Before going into the details of proving the renormalizability of quasi-PDFs, we first briefly explain the complexity of quasi-PDFs' renormalization:\\
{\bf 1) Lorentz symmetry is broken.} Because of the explicitly ``$z$"-direction dependence of quasi-PDFs, to identify all possible UV divergences, we need to study both four-dimensional loop momentum integration and three-dimensional loop momentum integration (with ``$z$"-direction fixed)\footnote{Note that these are the only two cases that we need to study. Other cases of loop momentum integration cannot generate new UV divergences.} for each individual loop. This amounts to $2^n$ different cases for a $n$-loop Feynman diagram, which is hard to handle.\\
{\bf 2) Renormalizaiton of composite operators is needed.} As an example, let's choose an $A_z=0$ axial gauge, and the quasi-quark PDFs become
\begin{align}\label{eq:ftqxAz}
\begin{split}
&\tilde{F}_{q/p}(\xi_z,\tilde{\mu}^2,p_z)= \frac{e^{i p_z \xi_z}}{p_z} \left. \langle h(p)| \overline{\psi}_q(\xi_z)\,\frac{\gamma_z}{2}\, \psi_q(0) | h(p) \rangle\right. .
\end{split}
\end{align}
For the renormalization of the quasi-quark PDFs, we need to renormalize the individual quark field, $\overline{\psi}_q(\xi_z)$ and $\psi_q(0)$ in Eq.~\eqref{eq:ftqxAz}, as well as the bi-local composite operator as a whole if it generates UV divergence.  The renormalization of the fields can be naturally taken care of by using the renormalized QCD Lagrangian in $A_z=0$ gauge.  It is the renormalization of the bi-local composite operators that is not covered by the renormalization of QCD.  This is effectively the same procedure for proving the renormalizability of PDFs, for example, for quark PDFs in $A^+=0$ gauge, one has to verify the multiplicative renormalization of the bi-local composite operators defining the quark-PDFs to prove their renormalizability.  It is the renormalization of the bi-local composite operators that mixes quark PDFs with gluon PDF \cite{Collins:1981uw}.  That is, the renormalizability of QCD Lagrangian itself is not enough to guarantee the renormalizability of quasi-PDFs, which are defined by composite operators.


In our explicit calculation and discussion in the rest of this paper, we choose the Feynman gauge since the renormalization of QCD Lagrangian in Feynman gauge is well known.  Since the renormalization of individual quark and gluon field is well defined, we will focus on the renormalization of the composite operators defining the quasi-PDFs.  We will mainly discuss coordinate-space quasi-PDFs, and thus whenever we mention quasi-PDFs,  we mean the coordinate-space quasi-PDFs.  We will first concentrate on quasi-quark PDFs, and then generalize our study to quasi-gluon PDF at the end.
